\documentclass[11pt]{article}
\usepackage[margin=0.85in]{geometry}
\usepackage{booktabs}
\usepackage{amsmath}
\usepackage[hidelinks]{hyperref}
\setlength{\parindent}{0pt}
\setlength{\parskip}{5pt}
\setlength{\tabcolsep}{7pt}

\begin{document}
\begin{center}
{\Large Wildfire sensitivity: working note}\\[4pt]
Until-stable run 3086411 $\cdot$ 18 September 2026 $\cdot$ two rounds, about 83 minutes
\end{center}

\textbf{Current result.}
The screening labels \textbf{LWP, LTS, and CO sensitive}, and leaves
\textbf{OMEGA undetermined}. Eight tested variables enter the insensitive set:
PBLH, CFLOW, CFMID, SWGDN, LWC, RH, T, and PREC. Together with the ten starting
assumptions (P, LAND, SLP, WS, TS, DUST, EPV, CHL, DMS, SO2EM), that gives
18 insensitive variables. The starting ten were not tested. These are screening
labels, not established causal effects. We use only until-stable experiments for
reported results.

\textbf{What is fitted and held out?}
For each candidate $c$, a global-latent LSTM engressor learns
$c=f(M,\varepsilon)$ on clean trajectories, where $M$ is the current insensitive
input set. Each trajectory contains 241 timesteps. We pool the original clean
training and test rows, then form four separate sets:

\begin{center}
\begin{tabular}{lrl}
\toprule
Clean set & Trajectories & Purpose \\
\midrule
Training & 23,564 & Fit the network \\
Validation & 3,110 & Early stopping; select best checkpoint \\
Recalibration & 1,350 & Fit the quantile correction \\
Control / test & 2,159 & Check corrected coverage \\
\bottomrule
\end{tabular}
\end{center}

The clean pool has 54,026 trajectories; the other 23,843 are unused after time-gap
purging and month matching. Recalibration and control are separately month-matched
to 6,786 wildfire trajectories. Splits use 30-day arrival-time blocks with history
gaps at both boundaries. Wildfire outcomes are used only for evaluation.

\textbf{What recalibration changes.}
Kuleshov, Fenner, and Ermon~\cite{kuleshov2018} fit a monotone map from predicted
CDF probabilities to observed frequencies, recommending isotonic regression.
Our implementation uses a simpler empirical quantile map. On the recalibration
rows, compute each observation's position $u_i=F_i(y_i)$ in the base model's
predictive distribution, approximated by randomized ranks among 400 draws.
Pool these positions across trajectories and timesteps, then use
\[
 p_q=\operatorname{Quantile}_q(\{u_i\}),
 \qquad Q_{\mathrm{corrected},x}(q)=Q_{\mathrm{base},x}(p_q).
\]
The network is unchanged. One mapping per candidate adjusts all its quantile
predictions. A central 90\% interval uses corrected 5th and 95th percentiles;
a 95\% interval uses corrected 2.5th and 97.5th percentiles. The same mapping is
applied to control and wildfire data. Fitting the correction does not guarantee
that it achieves the target coverage on the separate control set.

\textbf{Decision rule.}
A model is accepted if its control coverage at both 90\% and 95\% lies within
$\pm3$ estimated standard errors of nominal. Standard errors use 30-day blocks
and the same observation weighting as pooled coverage. A failed check leaves
$c$ undetermined and outside $M$. For accepted models, a control-minus-wildfire
95\% coverage drop greater than two percentage points is sensitive; otherwise
$c$ enters $M$. Sensitive and undetermined candidates are retried after $M$ grows.
The run stops when a round adds nobody.

\newpage
\textbf{Latest assessment of each tested candidate.}
Coverage is pooled across all 241 timesteps and reported in percent; drops are
percentage points. The eight insensitive decisions were made in round 1 and
retained. LWP, LTS, CO, and OMEGA were reassessed in round 2, which added nobody.
All numbers below come from this one until-stable run.

\begin{center}
\small
\begin{tabular}{lrrrrl}
\toprule
 & \multicolumn{2}{c}{Clean coverage} & Wildfire & & \\
Candidate & 90\% interval & 95\% interval & 95\% interval & Drop & Decision \\
\midrule
LWP   & 89.4 & 94.5 & 92.1 & $+2.3$ & Sensitive \\
LTS   & 88.1 & 94.2 & 91.3 & $+2.9$ & Sensitive \\
CO    & 81.0 & 87.4 & 75.7 & $+11.7$ & Sensitive \\
\midrule
PBLH  & 89.2 & 94.6 & 93.2 & $+1.4$ & Insensitive \\
CFLOW & 89.6 & 94.5 & 95.5 & $-1.0$ & Insensitive \\
CFMID & 91.2 & 95.2 & 94.7 & $+0.5$ & Insensitive \\
SWGDN & 88.7 & 94.3 & 94.6 & $-0.3$ & Insensitive \\
LWC   & 90.5 & 95.1 & 95.1 & $\approx0.0$ & Insensitive \\
RH    & 91.5 & 95.5 & 94.9 & $+0.6$ & Insensitive \\
T     & 85.5 & 92.4 & 92.0 & $+0.4$ & Insensitive \\
PREC  & 90.0 & 94.7 & 93.9 & $+0.9$ & Insensitive \\
\midrule
OMEGA & 92.7 & 96.5 & --- & --- & Undetermined \\
\bottomrule
\end{tabular}
\end{center}

\smallskip
Drops were calculated before rounding, so subtracting the displayed coverages
can differ by 0.1 point. Names omit the stored feature transforms: PBLH, LWP,
LWC, CO, PREC, DUST, and SO2EM use fifth roots; CHL and DMS use square roots.

\textbf{How much should we trust these labels?}
CO has the largest coverage drop, but its nominal 95\% interval covers only
87.4\% of clean control values. Its coverage varies enough between blocks that
the acceptance band is 85.0--105.0\%; it passes despite missing the target by
7.6 points. T similarly passes with 92.4\% coverage and a 90.9--99.1\% band.
An upper band above 100\% is an unbounded normal approximation, not a possible
coverage value. Passing means the test did not detect a departure, not that
coverage is necessarily close enough for our purpose.

OMEGA still overcovers: 92.7\% is outside its 88.3--91.7\% band, and 96.5\% is
outside 94.0--96.0\%. LWP and LTS are closer to their clean-coverage targets,
but their drops remain near the chosen two-point sensitivity threshold.

The next questions are whether CO and T need a stricter practical coverage
criterion, and whether the two-point drop is larger than clean-versus-clean
sampling variation. Also, if a smoke-affected variable enters $M$, it can help
predict another smoke-affected variable and hide that variable's shift.
Until-stable iteration does not remove this dependence on earlier decisions.
Month matching and pooled calibration support the comparison, but do not by
themselves isolate a causal wildfire effect.

\textbf{Files.}
\texttt{experiments/wildfire\_sensitivity/main.py} runs the screening.
The source log is
\nolinkurl{slurm/wildfire_sensitivity/run_until_stable_calibration_3086411.log}.
Checkpoints, mappings, and split indices are under\\
{\footnotesize\nolinkurl{/storage3/fs1/myu/Active/felixhu/weather_checkpoints/wildfire_sensitivity/3086411/}}.
To rerun: \texttt{sbatch slurm/wildfire\_sensitivity/run\_until\_stable.slurm}.

\begin{thebibliography}{1}
\small
\bibitem{kuleshov2018}
V.~Kuleshov, N.~Fenner, and S.~Ermon (2018).
\emph{Accurate Uncertainties for Deep Learning Using Calibrated Regression}.
ICML, PMLR 80, pp.~2796--2804.
\url{https://proceedings.mlr.press/v80/kuleshov18a.html}.
\end{thebibliography}
\end{document}
